\documentclass[10pt,a4paper]{article}

\usepackage{kotex}
\usepackage{fontspec}
\setmainfont{Times New Roman}
\setmainhangulfont{AppleMyungjo}
\setsanshangulfont{Apple SD Gothic Neo}

\usepackage[a4paper,top=20mm,bottom=20mm,left=20mm,right=20mm]{geometry}
\usepackage{array}
\usepackage{booktabs}
\usepackage{enumitem}
\linespread{1.15}
\setlength{\parindent}{0pt}

\newcommand{\sol}[3]{%
  \vspace{4pt}%
  {\bfseries\sffamily #1번} \quad 정답 \fbox{#2}\par
  \vspace{1pt}{\sffamily #3}\par}

\begin{document}

\begin{center}
{\fontsize{18}{20}\selectfont\bfseries\sffamily 2026 6월 고1 영어 변형 모의고사 — 정답 및 해설}\\[2pt]
{\sffamily (테스트본 18\,$\sim$\,22번)}
\end{center}
\vspace{6pt}

\begin{center}
\renewcommand{\arraystretch}{1.4}
\begin{tabular}{|c|c|c|c|c|c|}
\hline
\sffamily 문항 & 18 & 19 & 20 & 21 & 22 \\\hline
\sffamily 정답 & ④ & ② & ⑤ & ③ & ② \\\hline
\sffamily 유형 & \sffamily 목적 & \sffamily 심경 & \sffamily 주장 & \sffamily 함의 & \sffamily 요지 \\\hline
\sffamily 배점 & 2 & 2 & 2 & 3 & 2 \\\hline
\end{tabular}
\end{center}

\vspace{10pt}
\hrule
\vspace{8pt}

\sol{18}{④}{\textbf{글의 목적.} Greenfield Hiking Club 회장이 비로 손상된 등산로를 보수하는 ``trail repair day''에 회원들의 참여(sign up, spare a few hours)를 요청하는 편지다. 따라서 정답은 ④. 안전 수칙 안내(①)나 동호회 가입 권유(②)가 아니며, 폐쇄(③)·장비 기부(⑤)는 언급되지 않았다.}

\sol{19}{②}{\textbf{심경 변화.} 오디션 전 ``hands trembled, heart pounded, could hardly breathe''로 \textit{긴장(nervous)} 상태였다가, 연주를 성공적으로 마치고 ``walked off with a wide smile, glad that her effort had paid off''로 \textit{뿌듯함(proud)}을 느낀다. 따라서 ②. ‘nervous → proud’.}

\sol{20}{⑤}{\textbf{필자 주장.} 부모가 자녀의 모든 어려움을 대신 해결해 주면 배움의 기회를 빼앗는다며, 일상적인 문제(everyday troubles)는 자녀가 스스로 결과를 감당하도록 한 발 물러서라(stepping back)고 주장한다. 따라서 ⑤. 안전(①)은 ‘real danger’에 한정한 부차적 언급이다.}

\sol{21}{③}{\textbf{함의 추론.} ``carry our own weather''는 바로 뒤 ``we hold the thermostat of our inner climate''로 풀이된다. 외부 상황(storm)은 통제할 수 없지만 자신의 내면 상태는 스스로 정한다는 의미이므로 ③ ‘determine our inner state regardless of circumstances’. 환경 자체를 바꾼다는 ②와 혼동하지 않도록 주의.}

\sol{22}{②}{\textbf{글의 요지.} 빈 시간을 낭비로 여기지 말라는 글로, 아무것도 하지 않고 생각이 떠돌게 두면 뇌가 배경에서 연결을 만들어 좋은 아이디어(great ideas)가 떠오른다고 한다. 따라서 ② ‘충분한 여유 시간은 창의적인 생각이 떠오르는 바탕이 된다.’}

\end{document}
